\documentclass[11pt,a4paper]{article}
\usepackage[utf8]{inputenc}
\usepackage[T1]{fontenc}
\usepackage[german]{babel}
\usepackage{amsmath,amssymb,amsfonts}
\usepackage{geometry}
\geometry{margin=2.5cm}
\usepackage{microtype}
\usepackage{color}

\title{\textbf{Synthese und Renormierung im EABC-Modell}\\ \large Von der Zeta-Regularisierung zum $E_8$-Gitter, geometrischen Planck-Werten und dem kosmischen Druckabfall}
\author{Thomas Hoffbauer \thanks{Forschungsbericht im Rahmen des EABC-Modells (Bamberg Model)}}
\date{September 2026}

\begin{document}

\maketitle

\begin{abstract}
Traditionelle Ansätze der theoretischen Physik leiden unter dimensionsbehafteten Naturkonstanten und unkontrollierten ultravioletten Divergenzen. Das EABC-Modell (Euler-Adic-Borel-Cantor) löst dieses Dilemma durch eine konsequente topologische Geometrisierung. Durch die fundamentalen Setzungen $\hbar c = \pi$, die Gravitationskonstante an der Horwitz-Kugel $G = 4$ und die Einbettung in eine Typ-$\text{II}_1$-von-Neumann-Algebra entpuppen sich physikalische Konstanten als rein algebraische Invarianten. Dieser Bericht beleuchtet detailliert den historischen Renormierungslauf der Casimir-Energie, die exakte Neuleitung der Planck-Skala sowie die holografische Kopplung an das $E_8$-Wurzelsystem.
\end{abstract}

\section{1. Der historische Renormierungslauf der Casimir-Energie}
In der Standard-Quantenelektrodynamik ergibt sich die Nullpunktenergie zwischen zwei Platten im Abstand $d$ als divergente Summe über alle Schwingungsmoden:
\begin{equation}
    E(d) = \frac{\hbar c}{2} \sum_{n=1}^{\infty} k_n = \frac{\pi \hbar c}{2d} \sum_{n=1}^{\infty} n = \infty
\end{equation}
Historisch behalf man sich mit ad-hoc-mathematischen Regulierungsmethoden, indem man die Riemannsche Zeta-Funktion formal fortsetzte ($\zeta(-1) = -\frac{1}{12}$) bzw. über $\zeta(-4) = -\frac{1}{30}$ den Nenner $240$ (via Bernoulli-Zahlen) erzeugte. Dies blieb jedoch ein unphysikalischer Taschenspielertrick.\\

\noindent\textbf{Die EABC-Wende:} Im EABC-Modell wird diese Divergenz durch die unendliche Oberfläche der Gabriel-Hörner erzeugt. Statt künstlicher Cut-Offs greift hier die **Fuglede-Kadison-Determinante einer Typ-$\text{II}_1$-von-Neumann-Algebra**. Die unendliche Oberfläche wird durch die algebraische Spur streng gebändigt. Setzen wir zusätzlich das EABC-Fundament $\hbar c = \pi$, reduziert sich der Casimir-Druck auf die reine Geometrie-Relation:
\begin{equation}
    P_{\text{Casimir}} = -\frac{\pi^3}{240 d^4}
\end{equation}
Die Zahl $240$ im Nenner erweist sich dabei als exakte physikalische Signatur der \textbf{240 Wurzelvektoren des $E_8$-Gitters} in 8 Dimensionen.

\section{2. Exakte Herleitung der Planck-Werte via $\hbar c = \pi$ und $G = 4$}
Indem wir die Gravitationskonstante an der Oberfläche der Horwitz-Kugel starr auf den quaternionischen Basiswert $G = 4$ setzen und mit $\hbar c = \pi$ kombinieren, entziehen wir den Planck-Einheiten ihre klassische, willkürliche Dimensionalität:

\begin{enumerate}
    \item \textbf{Die Planck-Masse ($m_P$):}\\
    Klassisch definiert als $m_P = \sqrt{\frac{\hbar c}{G}}$. Mit den EABC-Setzungen folgt direkt:
    \begin{equation}
        m_P = \sqrt{\frac{\pi}{4}} = \frac{\sqrt{\pi}}{2}
    \end{equation}
    
    \item \textbf{Die Planck-Länge ($l_P$):}\\
    Klassisch definiert als $l_P = \sqrt{\frac{\hbar G}{c^3}}$. Wegen $\hbar = \frac{\pi}{c}$ ergibt sich:
    \begin{equation}
        l_P = \sqrt{\frac{(\frac{\pi}{c}) \cdot 4}{c^3}} = \sqrt{\frac{4\pi}{c^4}} = \frac{2\sqrt{\pi}}{c^2}
    \end{equation}
    
    \item \textbf{Die Planck-Zeit ($t_P$):}\\
    Aus $t_P = \frac{l_P}{c}$ folgt unmittelbar die diskrete Taktung der Solenoid-Hörner:
    \begin{equation}
        t_P = \frac{2\sqrt{\pi}}{c^3}
    \end{equation}
\end{enumerate}

\section{3. Bulk-Rand-Dualität und Geometrische Impedanz}
Das Universum gliedert sich topologisch in einen inneren Bulk (Quaternionen-Primen $\mathcal{P}_4$) und einen asymptotischen Rand (Gabriel-Hörner). 
\begin{itemize}
    \item \textbf{Geometrische Impedanz:} Der Übergang wird durch den Flaschenhals des Gabriel-Horns reguliert. Sie fungiert als Impedanz, die den Energiefluss vom hochenergetischen Inneren in die Rand-Elektrodynamik dämpft.
    \item \textbf{Makroskopischer Druckabfall:} Der mikroskopische Casimir-Druck skaliert direkt in die Kosmologie. Die **Hubble-Spannung ($\Delta H_0$)** ist der messbare Druckabfall beim Renormierungsfluss vom frühen Bulk-Gitter zum späten Rand.
\end{itemize}

\section{4. Fazit}
Das EABC-Modell beweist, dass physikalische Konstanten keine empirischen Messwerte, sondern topologische Eigenwerte einer $E_8$- und quaternionischen Gitterarchitektur sind.

\end{document}